\section{Background}

\paragraph{Mess3 process.} A 3-state nonunifilar HMM emitting $z \in \{0,1,2\}$, parameterized by emission clarity $\alpha$ and state persistence $x$, with $\beta=(1-\alpha)/2$, $y=1-2x$. The labeled transition matrices $T^{(z)}_{ij} = P(\text{emit } z, \text{go to } j \mid \text{state } i)$ are given in Appendix~\ref{app:mess3}. Key properties: nonunifilar (optimal prediction requires belief tracking), belief states form a Sierpinski gasket in the 2-simplex, and the spectral gap $\zeta = 1-3x$ controls belief convergence rate \cite{riechers2018spectral}.

\paragraph{Belief geometry in transformers.} Shai et al.~\cite{shai2024transformers} showed transformers linearly represent belief states $\veta^{(z_{1:t})} = \veta^{(\varnothing)} T^{(z_1)} \cdots T^{(z_t)} / \text{norm}$ in their residual stream. Piotrowski et al.~\cite{piotrowski2025constrained} showed this is a constrained Bayesian approximation: $r_1(z_{1:d}) = \bpi + \sum_{s=1}^{d} (\bpi T^{z_s} T^{d-s} - \bpi)$, where corrections decay as $\zeta^{d-s}$.

\paragraph{Factored World Hypothesis.} Transformers have an inductive bias toward factored representations in orthogonal subspaces \cite{shai2026factored}. For $N$ conditionally independent factors, the factored representation requires $\sum_n (d_n - 1)$ dimensions vs.\ $\prod_n d_n - 1$ (joint). This persists even when it sacrifices fidelity.

\paragraph{ICL as source discrimination.} Non-ergodicity forces hierarchical inference: identify the source, then predict within it. This is structurally identical to in-context learning. ICL and activation steering both modify the model's belief state, operating additively in log-belief space.
